\cotefile{63}
\scanpage{16}
\begin{center}
COURBES ELLIPTIQUES : FORMULAIRE

d'après J. Tate\\
mis au goût du jour par P. Deligne
\end{center}

1.\qquad \uline{Formules générales}.

Dans ce formulaire, on appellera courbe de genre $1$ sur un schéma de base $\mathrm{S}$ un morphisme propre et plat de présentation finie
\[
\mathrm{p}:\mathrm{E}\handwritten{\longrightarrow}\mathrm{S}\qquad,
\]
muni d'une section $\mathrm{e}$ contenue dans l'ouvert de lissité, dont les fibres géométriques sont des courbes réduites irréductibles de \handwritten{\uline}{genre arithmétique}\reversemarginpar\margin{\handwritten{?}}\normalmarginpar $1$. Les fibres géométriques de $\mathrm{p}$ sont donc soit :

a) une courbe elliptique, i.e. une courbe propre, lisse et connexe de genre $1$,

b) une droite projective dont on a identifié deux points distincts (cubique dans $\mathrm{P}_{2}$ à point double ordinaire),

c) une droite projective dont on a identifié deux points infiniment voisins (cubique dans $\mathrm{P}_{2}$ à singularité cuspidale).

On dira dans le cas a) (resp. b), c)) que la fibre est de type elliptique (resp. multiplicatif, resp. additif).

Si $\mathrm{E}$ est une courbe de genre $1$ sur $\mathrm{S}$, on définit un faisceau inversible $\omega$ sur $\mathrm{S}$ par la formule
\[
\omega=\mathrm{e}^{*}\,\Omega^{1}_{\mathrm{E}/\mathrm{S}}\qquad.
\]

La section $\mathrm{e}$ est un diviseur de Cartier relatif de $\mathrm{E}$ sur $\mathrm{S}$ ;
\clearpage
% archive sheet number: 15
\scanpage{17}\mspage{2}
de plus, on vérifie fibre par fibre que
\[
\mathrm{R}^{1}\,\mathrm{p}_{*}\,\handwritten{\mathcal{O}}(\mathrm{ne})=0\quad\text{pour}\quad\mathrm{n}>0\qquad.
\]

Il en résulte, par Riemann-Roch, que pour $\mathrm{n}>0$, le faisceau $\mathrm{p}_{*}\handwritten{\mathcal{O}}(\mathrm{ne})$ est looalement libre de rang $\mathrm{n}$, et que sa formation commute à tout changement de base. La suite exacte de cohomologie fournit de plus une suite exacte
\[
0\handwritten{\longrightarrow}\mathrm{p}_{*}\handwritten{\mathcal{O}}(\mathrm{ne})
\handwritten{\longrightarrow}\mathrm{p}_{*}\handwritten{\mathcal{O}}((\mathrm{n}+1)\mathrm{e})
\handwritten{\longrightarrow}\omega^{\handwritten{\otimes}-(\mathrm{n}+1)}
\handwritten{\longrightarrow}0\quad(\mathrm{n}>0)\,.
\]

Par ailleurs, $\handwritten{\mathcal{O}}(\mathrm{e})$ étant localement libre de rang $1$, on vérifie fibre par fibre, compte tenu de (SGA 1 X 1.2), que
\[
\handwritten{\mathcal{O}}_{\mathrm{S}}\handwritten{\xrightarrow{\sim}}
\mathrm{p}_{*}\handwritten{\mathcal{O}}_{\mathrm{E}}\handwritten{\xrightarrow{\sim}}
\mathrm{p}_{*}\handwritten{\mathcal{O}}(\mathrm{e})\qquad.
\]

En d'autres termes, les $\mathrm{p}_{*}\handwritten{\mathcal{O}}(\mathrm{me})$ pour $1\leq\mathrm{m}\leq\mathrm{n}$ définissent une filtration de $\mathrm{p}_{*}\handwritten{\mathcal{O}}(\mathrm{ne})$, de gradué associé
\makeatletter\tagsleft@true\makeatother
\[
\mathrm{Gr}\;\mathrm{p}_{*}\handwritten{\mathcal{O}}(\mathrm{ne})
=\sum_{\substack{\mathrm{i}=0\\\mathrm{i}\handwritten{\neq}1}}^{\mathrm{n}}\omega^{\handwritten{\otimes}-\mathrm{i}}\qquad.
\tag*{(1.1)}
\]

Le morphisme $\mathrm{p}$ est d'intersection complète relative ; il en résulte que $\mathrm{Rp}^{!}\handwritten{\mathcal{O}}_{\mathrm{S}}\in\mathrm{D}^{+}(\mathrm{E})$ a un seul groupe de cohomologie non\margin{\handwritten{\unsure{réf|vérif}}} nulle, $\mathrm{R}^{1}\,\mathrm{p}^{!}\handwritten{\mathcal{O}}_{\mathrm{S}}$, et que ce dernier est un faisceau inversible sur $\mathrm{E}$ ; on l'appelle le faisceau des différentielles régulières sur $\mathrm{E}$ et on le désigne par $\Omega^{\mathrm{reg}}_{\mathrm{E}/\mathrm{S}}$\margin{\handwritten{$\omega_{\mathrm{E}/\mathrm{S}}$}}. Là où $\mathrm{p}$ est lisse, on a $\Omega^{\mathrm{reg}}_{\mathrm{E}/\mathrm{S}}=\Omega^{1}_{\mathrm{E}/\mathrm{S}}$. La dualité montre que $\mathrm{R}^{1}\mathrm{p}_{*}\Omega^{\mathrm{rcg}}_{\mathrm{E}/\mathrm{S}}$ est isomorphe à $\handwritten{\mathcal{O}}$, $\mathrm{p}_{*}\Omega^{\mathrm{reg}}_{\mathrm{E}/\mathrm{S}}$ est donc \handwritten{/} localement libre, de formation compatible à tout changement\margin{\handwritten{\unsure{réf|ref}}} de base et, raisonnant fibre par fibre, on en déduit que
\clearpage
% archive sheet number: 16
\scanpage{18}\mspage{3}
\[
\mathrm{p}_{*}\,\Omega^{\mathrm{reg}}_{\mathrm{E}/\mathrm{S}}\handwritten{\xrightarrow{\sim}}\mathrm{o}^{*}\,\Omega^{1}_{\mathrm{E}/\mathrm{S}}=\omega\quad,
\]
et que pour la flèche adjointe à la flèche inverse
\[
\mathrm{p}^{*}\,\omega\handwritten{\xrightarrow{\sim}}\Omega^{\mathrm{reg}}_{\mathrm{E}/\mathrm{S}}\qquad.
\]
\margin{\handwritten{?}}
Les sections de $\Omega^{\mathrm{reg}}_{\mathrm{E}/\mathrm{S}}$ ainsi associées aux sections de $\omega$ s'appellent les différentielles invariantes.

Soit $\unsure{\grave{\pi}|\dot{\pi}|\pi}$ une section inversible de $\omega$ ; il existe, localement pour la topologie de Zariski, une base $(\mathrm{y},\mathrm{x},1)$ de $\mathrm{p}_{*}\handwritten{\mathcal{O}}(3\mathrm{e})$ telle que l'image de $\mathrm{y}$ dans $\omega^{\handwritten{\otimes}-3}$ soit $-\pi^{\handwritten{\otimes}-3}$, que $\mathrm{x}$ appartienne à $\mathrm{p}_{*}\handwritten{\mathcal{O}}(2\mathrm{e})$, et que son image dans $\omega^{\handwritten{\otimes}-2}$ soit $\pi^{\handwritten{\unsure{\otimes|\circ}}-2}$. Si $\pi'$ est une autre section inversible de $\omega$, les bases correspondantes s'écrivent toutes, sous forme implicite (pour $\mathrm{r},\mathrm{s},\mathrm{t}$ convenables)
\[
\handwritten{\left\{}\begin{aligned}
\mathrm{x}&=\mathrm{u}^{2}\,\mathrm{x}'+\mathrm{r}\\
\mathrm{y}&=\mathrm{u}^{3}\,\mathrm{y}'+\mathrm{s}\,\mathrm{u}^{2}\,\mathrm{x}'+\mathrm{t}\\
\pi'&=\mathrm{u}\,\pi\qquad.
\end{aligned}\right.
\tag*{(1.2)}
\]

Le faisceau\handwritten{$\backslash$} $\handwritten{\mathcal{O}}(3\mathrm{e})$ est très ample et $(\mathrm{x},\mathrm{y},\tikz[baseline=(digit.base)]{\node[inner sep=0pt] (digit) {$1)$};\handwritten{\draw[thin,gray,rotate around={-16:(digit.center)}] (digit.center) ellipse [x radius=.68em,y radius=1.35em];}}$ sont les coordonnées projectives d'un plongement de $\mathrm{E}$ dans $\mathrm{P}_{2}/\mathrm{S}$.

La courbe image vérifiera une équation
\[
\mathrm{y}^{2}+\mathrm{a}_{1}\,\mathrm{x}\,\mathrm{y}+\mathrm{a}_{3}\,\mathrm{y}
=\mathrm{x}^{3}+\mathrm{a}_{2}\,\mathrm{x}^{2}+\mathrm{a}_{4}\,\mathrm{x}+\mathrm{a}_{6}\qquad.
\tag*{(1.3)}
\]

Pour vérifier l'existence et l'unicité de cette équation, on
\clearpage
% archive sheet number: 17
\scanpage{19}\mspage{4}
utilise (1.1), en regardant les deux membres de (1.3) comme des sections de $\mathrm{p}_{*}\handwritten{\mathcal{O}}(6\mathrm{e})$ : en effet, $\mathrm{y}^{2}-\mathrm{x}^{3}$ est dans $\mathrm{p}_{*}\handwritten{\mathcal{O}}(5\mathrm{e})$, $\mathrm{xy}$ dans $\mathrm{p}_{*}\handwritten{\mathcal{O}}(5\mathrm{e})$, $\mathrm{x}^{2}$ dans $\mathrm{p}_{*}\handwritten{\mathcal{O}}(4\mathrm{e})$, $\mathrm{y}$ dans $\mathrm{p}_{*}\handwritten{\mathcal{O}}(3\mathrm{e})$, $\mathrm{x}$ dans $\mathrm{p}_{*}\handwritten{\mathcal{O}}(2\mathrm{e})$ et $1$ dans $\mathrm{p}_{*}\handwritten{\mathcal{O}}$.

Si on pose
\[
\mathrm{F}(\mathrm{x},\mathrm{y})=\mathrm{y}^{2}+\mathrm{a}_{1}\,\mathrm{xy}+\mathrm{a}_{3}\,\mathrm{y}-\mathrm{x}^{3}-\mathrm{a}_{2}\,\mathrm{x}^{2}-\mathrm{a}_{4}\,\mathrm{x}-\mathrm{a}_{6}
\]
la forme différentielle invariante définie par $\pi$ s'écrit
\[
\begin{aligned}
\pi&=\mathrm{dx}/\mathrm{F}_{\mathrm{y}}=-\mathrm{dy}/\mathrm{F}_{\mathrm{x}}\qquad\text{soit}\\[.65em]
\pi&=\handwritten{\frac}{\mathrm{dx}}{2\mathrm{y}+\mathrm{a}_{1}\mathrm{x}+\mathrm{a}_{3}}
=\handwritten{\frac}{\mathrm{dy}}{3\mathrm{x}^{2}+2\mathrm{a}_{2}\mathrm{x}+\mathrm{a}_{4}-\mathrm{a}_{1}\mathrm{y}}\qquad.
\end{aligned}
\]

Réciproquement toute équation (1.3) est l'équation non homogène d'une courbe de genre $1$ plongée comme une cubique dans $\mathrm{P}_{2}$ et présentant un point d'inflexion en $(0,1,0)$\margin{\handwritten{$\Big)$\,?}} avec la droite de l'infini pour tangente d'inflexion.

On pose alors
% A short isolated stroke at left faces the row beginning d_6 below.
\reversemarginpar
\[
\mathrm{b}_{2}=\mathrm{a}_{1}^{2}+4\mathrm{a}_{2}\qquad
\mathrm{b}_{4}=\mathrm{a}_{1}\mathrm{a}_{3}+2\mathrm{a}_{4}\qquad
\mathrm{b}_{6}=\mathrm{a}_{3}^{2}+4\mathrm{a}_{6}
\tag*{(1.4)}
\]
\[
\begin{aligned}
&-\mathrm{b}_{8}=\mathrm{a}_{1}\,\mathrm{a}_{3}\,\mathrm{a}_{4}+\mathrm{a}_{4}^{2}-\mathrm{a}_{1}^{2}\,\mathrm{a}_{6}-\mathrm{a}_{2}\,\mathrm{a}_{3}^{2}-\overset{\inserted{\handwritten{2^{2}}}}{4}\mathrm{a}_{2}\,\mathrm{a}_{6}\\[.55em]
&\mathrm{d}_{6}=\mathrm{b}_{2}\,\mathrm{b}_{4}-\overset{\inserted{\handwritten{2^{1}.3^{2}}}}{18}\mathrm{b}_{6}\qquad
\mathrm{d}_{8}=\mathrm{b}_{2}^{2}\,\mathrm{b}_{4}-2^{4}.3\,\mathrm{b}_{4}^{2}+2.3^{2}\,\mathrm{b}_{2}\,\mathrm{b}_{6}
\end{aligned}
\]
\margin{\handwritten{$\backslash$}}\normalmarginpar

On a
\[
\begin{aligned}
&-4\mathrm{b}_{8}=\mathrm{b}_{4}^{2}-\mathrm{b}_{2}\,\mathrm{b}_{6}\\[.45em]
&\tikz[baseline=(digit.base)]{\node[inner sep=0pt] (digit) {$1$};\handwritten{\draw[thin,gray,rotate around={-12:(digit.center)}] (digit.center) ellipse [x radius=.43em,y radius=1.32em];}}2\,\mathrm{d}_{6}=\mathrm{c}_{6}+\mathrm{b}_{2}\,\mathrm{c}_{4}\qquad
-\tikz[baseline=(digit.base)]{\node[inner sep=0pt] (digit) {$1$};\handwritten{\draw[thin,gray,rotate around={-12:(digit.center)}] (digit.center) ellipse [x radius=.4em,y radius=1.13em];}}2\,\mathrm{d}_{8}=\mathrm{c}_{4}^{2}+\mathrm{b}_{2}\,\mathrm{c}_{6}
\end{aligned}
\]
\clearpage
% archive sheet number: 18
\scanpage{20}\mspage{5}
\[
\mathrm{c}_{4}=\mathrm{b}_{2}^{2}-\overset{\inserted{\handwritten{2^{3}3}}}{24}\,\mathrm{b}_{4}\qquad
-\mathrm{c}_{6}=\mathrm{b}_{2}^{3}-\overset{\inserted{\handwritten{2^{2}3^{2}}}}{36}\,\mathrm{b}_{2}\,\mathrm{b}_{4}
+\overset{\inserted{\handwritten{2^{3}3^{3}}}}{\tikz[baseline=(digits.base)]{\node[inner sep=0pt] (digits) {$216$};\handwritten{\draw[thin,gray,rotate around={-34:(digits.center)}] (digits.center) ellipse [x radius=1.24em,y radius=1.65em];}}}\mathrm{b}_{6}
\tag*{(1.5)}
\]
\[
\begin{aligned}
\Delta&=-\mathrm{b}_{2}^{2}\,\mathrm{b}_{8}-\overset{\inserted{\handwritten{2^{3}}}}{8}\,\mathrm{b}_{4}^{3}
-\overset{\inserted{\handwritten{3^{3}}}}{27}\,\mathrm{b}_{6}^{2}+\overset{\inserted{\handwritten{3^{2}}}}{9}\,\mathrm{b}_{2}\,\mathrm{b}_{4}\,\mathrm{b}_{6}\\[.6em]
&=\mathrm{b}_{4}^{3}-\overset{\inserted{\handwritten{3^{3}}}}{27}\,\mathrm{b}_{6}^{2}
+\mathrm{b}_{8}\bigl(\overset{\inserted{\handwritten{2^{2}3^{2}}}}{36}\,\mathrm{b}_{4}-\mathrm{b}_{2}^{\tikz[baseline=(digit.base)]{\node[inner sep=0pt] (digit) {$\scriptstyle\handwritten{2}$};\handwritten{\draw[thin,gray,rotate around={-8:(digit.center)}] ([yshift=.36em]digit.center) ellipse [x radius=.52em,y radius=.84em];}}}\bigr)\quad\handwritten{\raisebox{.7em}{?}}\quad.
\end{aligned}
\]

Soient maintenant $(\mathrm{x}',\mathrm{y}',\pi')$ reliés à $(\mathrm{x},\mathrm{y},\pi)$ par (1.2). On aura :
\[
\begin{aligned}[t]
\mathrm{u}\,\mathrm{a}'_{1}&=\mathrm{a}_{1}+2\mathrm{s}\\
\mathrm{u}^{2}\,\mathrm{a}'_{2}&=\mathrm{a}_{2}-\mathrm{sa}_{1}+3\mathrm{r}-\mathrm{s}^{2}\\
\mathrm{u}^{3}\,\mathrm{a}'_{3}&=\mathrm{a}_{3}+\mathrm{ra}_{1}+2\mathrm{t}=\mathrm{F}_{\mathrm{t}}(\mathrm{r},\mathrm{t})\\
\mathrm{u}^{4}\,\mathrm{a}'_{4}&=\mathrm{a}_{4}-\mathrm{sa}_{3}+2\mathrm{a}_{2}\mathrm{r}-(\mathrm{t}+\mathrm{rs})\unsure{\acute{\mathrm{a}}|{}'\mathrm{a}|\dot{\mathrm{a}}|\mathrm{a}}_{1}+3\mathrm{r}^{2}-2\mathrm{st}
=-\mathrm{F}(\mathrm{r},\mathrm{t})-\mathrm{s}\,\mathrm{F}_{\mathrm{t}}(\mathrm{r},\mathrm{t})\\
\mathrm{u}^{6}\,\mathrm{a}'_{6}&=\mathrm{a}_{6}+\mathrm{ra}_{4}+\mathrm{r}^{2}\mathrm{a}_{2}+\mathrm{r}^{3}-\mathrm{ta}_{3}-\mathrm{t}^{2}-\mathrm{rta}_{1}=-\mathrm{F}(\mathrm{r},\mathrm{t})
\end{aligned}
\tag*{(1.6)}
\]
\[
\begin{aligned}[t]
\mathrm{u}^{2}\,\mathrm{b}'_{2}&=\mathrm{b}_{2}+12\,\mathrm{r}\\
\mathrm{u}^{4}\,\mathrm{b}'_{4}&=\mathrm{b}_{4}+\mathrm{rb}_{2}+6\mathrm{r}^{2}\quad(\text{équ. de discr. }\mathrm{c}_{4})\\
\mathrm{u}^{6}\,\mathrm{b}'_{6}&=\mathrm{b}_{6}+2\mathrm{r}\,\mathrm{b}_{4}+\mathrm{r}^{2}\,\mathrm{b}_{2}+4\mathrm{r}^{3}\quad(\text{équ. de discr. }16\,\Delta)\\
\mathrm{u}^{8}\,\mathrm{b}'_{8}&=\mathrm{b}_{8}+3\mathrm{r}\,\mathrm{b}_{6}+3\mathrm{r}^{2}\,\mathrm{b}_{4}+\mathrm{r}^{3}\,\mathrm{b}_{2}+3\mathrm{r}^{4}\\
\mathrm{u}^{6}\,\mathrm{d}'_{6}&=\mathrm{d}_{6}+\mathrm{c}_{4}\,\mathrm{r}\qquad\mathrm{u}^{8}\,\mathrm{d}'_{8}=\mathrm{d}_{8}-\mathrm{c}_{6}\,\mathrm{r}
\end{aligned}
\tag*{(1.7)}
\]
\clearpage
% archive sheet number: 13
% The second figure also admits 9; see the doubts table.
\scanpage{21}\mspage{6}
\[
\mathrm{u}^{4}\,\mathrm{c}'_{4}=\mathrm{c}_{4}\qquad
\mathrm{u}^{6}\,\mathrm{c}'_{6}=\mathrm{c}_{6}\qquad
\mathrm{u}^{12}\,\Delta'=\Delta\qquad.
\tag*{(1.8)}
\]

Les formules (1.8) expriment que la section $\mathrm{c}_{4}\,\pi^{\handwritten{\otimes}4}$ (resp. $\mathrm{c}_{6}\,\pi^{\handwritten{\otimes}6}$, resp. $\Delta\,\pi^{\handwritten{\otimes}12}$) du fibré $\omega^{\handwritten{\unsure{\otimes|\text{\textcircled{\scriptsize$\backslash$}}|\circ}}4}$ (resp. $\omega^{\handwritten{\otimes}6}$, $\omega^{\handwritten{\otimes}12}$)\raisebox{.35ex}{\tiny.} ne dépend pas du choix arbitraire de $\pi$, $\mathrm{x}$ et $\mathrm{y}$.
\par\noindent\hfill
\margin{\hspace{1.8cm}\makebox[0pt][r]{\begin{minipage}[t]{10.1cm}
\diagram{\handwritten{\raisebox{2.65\baselineskip}{\uline{Sur un \unsure{corps}}}\ \(\left\{\begin{array}{@{}l@{\quad}l@{}}
\text{\unsure{cas} elliptique :}&\Delta\neq0\\
\text{cas multiplicatif :}&\Delta=0,\ \unsure{C|c}_{4}\neq0\ \text{(équiv.}\\[-.2ex]
&\hfill{\scriptstyle\unsure{c|e}}\\[-.3ex]
&\hspace{4em}\text{à }C_{6}\neq0)\\
\text{cas \unsure{additif} :}&\Delta=0,\ C_{4}=0\ \text{i.e.}\\
&\Delta=0,\ C_{6}=0
\end{array}\right.\)}}
\end{minipage}}}
\par\vspace{7.7\baselineskip}

2.\qquad \uline{Cas où 2 et 3 sont inversibles}.\margin{\handwritten{$\backslash$}}

Lorsque $2$ et $3$ sont inversibles, on peut d'une et d'une seule façon choisir $\mathrm{x}$ et $\mathrm{y}$ ($\pi$ étant donné) de sorte que dans l'équation (1.3.) on ait $\mathrm{a}_{1}=\mathrm{a}_{2}=\mathrm{a}_{3}=0$. On a alors
\[
\mathrm{y}^{2}=\mathrm{x}^{3}+\mathrm{a}_{4}\,\mathrm{x}+\mathrm{a}_{6}\qquad,\ \text{avec}\quad\pi=\mathrm{dx}/2\mathrm{y}\qquad,
\tag*{(2.1)}
\]
\[
\mathrm{c}_{4}=-2^{4}.3\mathrm{a}_{4}\quad,\quad
\mathrm{c}_{6}=-2^{5}.3^{3}\,\mathrm{a}_{6}\quad,\quad
\Delta=2^{4}(-4\mathrm{a}_{4}^{3}-27\,\mathrm{a}_{6}^{2})\quad.
\tag*{(2.2)}
\]

Posant $\mathrm{Y}=2\mathrm{y}$, $\mathrm{X}=\mathrm{x}$, on retombe sur la forme de Weierstrass
\[
\mathrm{Y}^{2}=4\,\mathrm{X}^{3}-\mathrm{g}_{2}\,\mathrm{X}-\mathrm{g}_{3}\quad,\ \text{avec}\quad\pi=\mathrm{d}\,\mathrm{X}/\mathrm{Y}\quad,\quad\text{et}
\tag*{(2.3)}
\]
\[
\begin{aligned}[t]
&\mathrm{g}_{2}=-4\,\mathrm{a}_{4}=\handwritten{\frac}{1}{12}\,\mathrm{c}_{4}\qquad
\mathrm{g}_{3}=-4\,\mathrm{a}_{6}=\handwritten{\frac}{1}{216}\,\mathrm{c}_{6}\\
&\Delta=\mathrm{g}_{2}^{3}-27\,\mathrm{g}_{3}^{2}\qquad.
\end{aligned}
\tag*{(2.4)}
\]

\uline{Proposition} 2.5.\quad \uline{Au dessus de} $\operatorname{Spec}(\handwritten{\mathbb{Z}}[2^{-1},\,3^{-1}]$, \uline{il existe un schéma de modules pour les courbes de genre un munies d'une forme différentielle invariante inversible}. \uline{\unsure{C;|Ce|C,} schéma est}
\[
\operatorname{Spec}(\handwritten{\mathbb{Z}}[2^{-1},\,3^{-1}][\mathrm{g}_{2},\,\mathrm{g}_{3}]\qquad,
\]
\clearpage
% archive sheet number: 20
\scanpage{22}\mspage{7}
\uline{avec pour courbe universelle} (\uline{en coordonnées non \unsure{homogò\`nos|homogònos|homogènes}})
\[
\mathrm{Y}^{2}=4\,\mathrm{X}^{3}-\mathrm{g}_{2}\,\mathrm{X}-\mathrm{g}_{3}
\]
\margin{\handwritten{\uline{\unsure{section} \unsure{unité}} ?}}
\uline{et pour différentielle invariante} $\pi=\mathrm{d}\,\mathrm{X}/\mathrm{Y}$.

De plus, lorsque $2$ et $3$ sont inversibles, la donnée d'une forme différentielle invariante suffit à rigidifier une courbe de genre $1$.

Il résulte aussitôt de (2.2) qu'on a
\[
\mathrm{c}_{4}^{3}-\mathrm{c}_{6}^{2}=\underset{\inserted{\handwritten{12^{3}}}}{1728}\,\Delta
\tag*{(2.6)}
\]
et cette formule, se réduisant à une identité algébrique, est valable en toute caractéristique.

3.\qquad \uline{Cas où} $2$ \uline{et} $\mathrm{c}_{4}$ \uline{sont inversibles}.

Supposons que $\mathrm{E}$ soit une courbe de genre un sur une base $\mathrm{S}$ où $2$ soit inversible. La forme différentielle invariante $\pi$ étant donnée, on peut choisir $\mathrm{x}$ et $\mathrm{y}$ tels que dans (1.3) on ait $\mathrm{a}_{1}=\mathrm{a}_{3}=0$. Les changements de coordonnées respectant cette condition s'obtiennent en faisant $\mathrm{u}=1$, $\mathrm{s}=\mathrm{t}=0$ dans (1.2).

On a donc
\[
\mathrm{y}^{2}=\mathrm{x}^{3}+\mathrm{a}_{2}\,\mathrm{x}^{2}+\mathrm{a}_{4}\,\mathrm{x}+\mathrm{a}_{6}\,\raisebox{.14ex}{\scriptsize\unsure{,|'}}
\tag*{(3.1)}
\]
\[
\mathrm{b}_{2}=4\mathrm{a}_{2}\qquad\mathrm{b}_{4}=2\mathrm{a}_{4}\qquad
\mathrm{b}_{6}=4\mathrm{a}_{6}\qquad\mathrm{d}_{6}=8(\mathrm{a}_{2}\,\mathrm{a}_{4}-9\,\mathrm{a}_{6})
\]
\[
\mathrm{c}_{4}=2^{4}(\mathrm{a}_{2}^{2}-3\mathrm{a}_{4})\qquad
-\mathrm{c}_{6}=2^{6}\,\mathrm{a}_{2}^{3}-2^{5}\,3^{2}\,\mathrm{a}_{2}\,\mathrm{a}_{4}+2^{5}\,3^{3}\,\mathrm{a}_{6}\quad\text{et}
\tag*{(3.2)}
\]
\[
\Delta=2^{4}\,\delta\qquad,
\]
\clearpage
% archive sheet number: 21
\scanpage{23}\mspage{8}
où $\delta$ est le discriminant du second membre de (3.1). Le discriminant $\delta$ étant invariant par translation, il suffit de vérifier cette identité algébrique pour $3$ inversible et $\mathrm{a}_{2}=0$.

Lorsque $\mathrm{c}_{4}$ est inversible, un et un seul choix de $\mathrm{x}$ et $\mathrm{y}$ donne $\mathrm{d}_{6}=0$ ? i.e. $\mathrm{a}_{2}\,\mathrm{a}_{4}-9\,\mathrm{a}_{6}=0$. On a, lorsque cette condition est remplie,
\[
-\mathrm{c}_{6}=2^{6}(\mathrm{a}_{2}^{3}-3^{3}\,\mathrm{a}_{6})\qquad
\Delta=2^{6}(-\mathrm{a}_{2}^{3}\,\mathrm{a}_{6}-\mathrm{a}_{4}^{3}+2.3^{3}\,\mathrm{a}_{6}^{2})
\tag*{(3.3)}
\]
et
\[
\begin{aligned}[t]
\mathrm{a}_{2}\,\mathrm{c}_{4}&=-2^{-2}\,\mathrm{c}_{6}\\
\mathrm{a}_{4}\,\mathrm{c}_{4}^{2}&=-2^{2}\,3^{2}\,\Delta\\
\mathrm{a}_{6}\,\mathrm{c}_{4}^{3}&=\mathrm{c}_{6}\,\Delta\qquad.
\end{aligned}
\tag*{(3.4)}
\]

Il suffit de vérifier ces identités en caractéristique $0$, auquel cas la troisième est le produit des deux premières, puisque $\mathrm{d}_{6}=0$.

\uline{Proposition} 3.5.\quad \uline{Au dessus de} $\operatorname{Spec}(\handwritten{\mathbb{Z}}[2^{-1}])$, \uline{il existe un schéma de modul}\tikz[baseline=(ending.base)]{\node[inner sep=0pt] (ending) {\uline{e\handwritten{s}}};\handwritten{\draw[overlay,rotate around={-6:(ending.center)}] (ending.center) ellipse [x radius=.6em,y radius=1.35em];}}\uline{ pour les courbes de genre un munies d'une forme différentielle invariante inversible et telles que} $\mathrm{c}_{4}$ \uline{soit inversible. Ce schéma} \reversemarginpar\margin{\handwritten{\tikz[baseline=-.15ex]{\draw[thin,gray] (0,0) -- (.63em,.14em);}}}\normalmarginpar\uline{est}
\[
\operatorname{Spec}(\handwritten{\mathbb{Z}}[2^{-1}][\mathrm{a}_{2},\,\mathrm{a}_{4},\,\mathrm{a}_{6},\,(\mathrm{a}_{2}^{2}-3\,\mathrm{a}_{4})^{-1}]/(\mathrm{a}_{2}\,\mathrm{a}_{4}-9\,\mathrm{a}_{6}))
\]
\uline{avec pour courbe universelle} (\uline{en coordonnées non homogè}\tikz[baseline=(ending.base)]{\node[inner sep=0pt] (ending) {\uline{ne}\handwritten{s})};\handwritten{\draw[overlay] ([xshift=1.3em,yshift=2.16em]ending.center) .. controls +(-.8em,.42em) and +(.3em,1.1em) .. ([xshift=-.63em,yshift=.2em]ending.center) .. controls +(-.72em,-1.6em) and +(-.76em,-1.6em) .. ([xshift=1.35em,yshift=-.65em]ending.center) .. controls +(.7em,.8em) and +(.5em,1.6em) .. ([xshift=1.32em,yshift=2.22em]ending.center);}}
\[
\mathrm{y}^{2}=\mathrm{x}^{3}+\mathrm{a}_{2}\,\mathrm{x}^{2}+\mathrm{a}_{4}\,\mathrm{x}+\mathrm{a}_{6}
\]
\clearpage
% archive sheet number: 22
\scanpage{24}\mspage{9}
\uline{et pour différentielle invariante} $\pi=\mathrm{dx}/2\mathrm{y}$.

De plus, sous ces hypothèses, la donnée d'une forme différentielle invariante suffit à rigidifier une courbe.

4.\qquad \uline{Cas où} $2$ \uline{est nilpotent et} $\mathrm{c}_{4}$ \uline{inversible}.

Lorsque $2=0$, la forme différentielle ne suffit plus jamais à rigidifier une courbe de genre un. Lorsque $\mathrm{c}_{4}$ est inversible et $2$ nilpotent, donc $\mathrm{a}_{1}$ inversible, en changeant le choix de $\pi$, $\mathrm{x}$, $\mathrm{y}$ on peut trouver une équation
\[
\mathrm{y}^{2}+\mathrm{xy}=\mathrm{x}^{3}+\mathrm{a}_{2}\,\mathrm{x}^{2}+\mathrm{a}_{6}\qquad
\pi=\handwritten{\frac}{\mathrm{dx}}{\mathrm{x}+2\mathrm{y}}\qquad.
\tag*{(4.1)}
\]

Les autres choix de $\mathrm{x}$ et $\mathrm{y}$ donnant lieu à une telle équation s'obtiennent en faisant, dans (1.2)
\[
\mathrm{r}=\mathrm{t}=0\qquad\mathrm{u}=1+2\mathrm{s}
\tag*{(4.2)}
\]
et on aura
\[
\begin{aligned}[t]
\mathrm{u}^{2}\,\mathrm{a}'_{2}&=\mathrm{a}_{2}-\mathrm{s}-\mathrm{s}^{2}\\
\mathrm{u}^{6}\,\mathrm{a}'_{6}&=\mathrm{a}_{6}
\end{aligned}
\tag*{(4.3)}
\]
et \uline{modulo 2}
\[
\mathrm{c}_{4}=1\qquad\mathrm{c}_{6}=1\qquad\Delta=\mathrm{a}_{6}\qquad.
\tag*{(4.4)}
\]
\margin{\handwritten{\unsure{T|$\overline{\mathrm{I}}$}\,\unsure{,|/}\,$\sim$}}

5.\qquad \uline{Cas où} $\mathrm{c}_{4}=0$, $\mathrm{p}=2$ \uline{ou} $\mathrm{p}=3$\tikz[baseline=(mark.base)]{\node[inner sep=0pt] (mark) {.};\handwritten{\draw[overlay,thin,gray] ([xshift=.1em,yshift=.86em]mark.center) .. controls +(-.2em,.5em) and +(-.83em,-.8em) .. ([xshift=-.23em,yshift=-.9em]mark.center) .. controls +(.67em,-.12em) and +(0,-.45em) .. ([xshift=.93em,yshift=.5em]mark.center);}}\margin{\handwritten{\tikz[baseline]{\draw[thin,gray] (0,0) .. controls (.12em,.38em) and (.38em,.27em) .. (.58em,.25em);}}}

a) $\mathrm{p}=3$

En caractéristique $3$, lorsque l'équation est mise sou\tikz[baseline=(ending.base)]{\node[inner sep=0pt] (ending) {s};\node[overlay,inner sep=0pt,anchor=base] at (ending.base) {\deleted{\unsure{r|i}}};\handwritten{\draw[overlay,thin,gray] ([xshift=-.1em,yshift=.85em]ending.center) .. controls +(-.35em,.32em) and +(-.45em,.82em) .. ([xshift=-.67em,yshift=-.05em]ending.center) .. controls +(.2em,-1em) and +(.8em,-.48em) .. ([xshift=.59em,yshift=-.1em]ending.center) .. controls +(0,.75em) and +(.42em,.34em) .. ([xshift=-.1em,yshift=.55em]ending.center);}} forme
% archive sheet number: 23
